\section{Experiments}
\subsection{Experimental Setup}
\textbf{Datasets.}
To validate the observations mentioned above and evaluate multi-foundation-model fusion across different fusion strategies, we conduct experiments on two TCGA datasets, including TCGA-UVM \citep{uvmdataset} and TCGA-BLCA \citep{blcatp53mutation}.
To examine whether above insights hold across diverse foundation-model compositions, we use TCGA-UVM as a diagnostic study setting and extract representations from 13 pathology foundation models.
This broad model coverage enables us to systematically evaluate how representation relationships and encoder composition relate to downstream fusion behavior.
Following previous molecular characterization studies, we formulate the task as D3/M3 subtype binary classification \citep{fouruvm}, with 38 D3 cases and 42 M3 cases.
We additionally evaluate CoR-Fuse on TCGA-BLCA, a larger dataset with 457 slides.
The task is formulated as binary prediction of TP53 mutation status from WSI representations, distinguishing TP53-mutant and TP53-wildtype cases \citep{blcatp53mutation}, with 238 mutant and 219 wildtype cases.

\textbf{Implementation Details.}
For both TCGA-UVM and TCGA-BLCA datasets, we perform patient-level stratified five-fold cross-validation and report the average performance across five folds. 
All foundation-model embeddings are extracted using the same preprocessing pipeline.
Performance is evaluated using the area under the receiver operating characteristic curve (AUC) and the area under the precision-recall curve (AUPRC). Computational efficiency is assessed in terms of trainable parameters and floating-point operations (FLOPs).
Detailed training configurations and hyperparameter settings are provided in Appendix \ref{implemented details}.
\subsection{Experimental Results}
\subsubsection{Multi-Foundation-Model Fusion under Different Representation Compositions}
\textbf{Individual Foundation Model Performance.}
We first evaluate all 13 FMs on TCGA-UVM (Appendix Table~\ref{tab:uvm_all_fm_performance}).
The individual performance varies substantially across foundation models, suggesting that individual FMs provide varying levels of task relevance and motivating representation-aware fusion.

\textbf{Plug-and-Play Compatibility.}
To evaluate the plug-and-play capability of CoR-Fuse, we examine its compatibility with both different foundation-model compositions and downstream prediction heads. 
On the downstream side, we combine CoR-Fuse with four downstream heads, including ABMIL, TransMIL \citep{transMIL}, MLP, and GNN-based classifiers.
The results in Appendix Table~\ref{tab:downstream_heads} show that CoR-Fuse can be effectively combined with different architectures.
On the encoder side, CoR-Fuse can adapt to different numbers and combinations of foundation models without requiring a fixed encoder set.

\textbf{Comparison with Existing Fusion Strategies.}
Based on the above analysis, we hypothesize that fusion effectiveness depends on the composition of integrated representations rather than the number of encoders alone.
To examine this hypothesis, we evaluate representative fusion strategies across different foundation-model compositions. 
The compared representative fusion paradigms (detailed in Appendix~\ref{sec:contrast_method}), include: \textbf{Fixed aggregation:} Mean fusion, Concatenation; \textbf{Global learned weighting:} Global weighting fusion; \textbf{Slide-specific weighting:} Gated fusion; \textbf{Interaction-based fusion:} Cross-attention, Self-attention.

We further compare CoR-Fuse with recently proposed multi-foundation-model fusion methods such as Meta-Encoder~\citep{gao2026metaencoder} and FuseCPath~\citep{yang2025fusecpath}.
Results in Appendix Table~\ref{tab:uvm_existing_methods} show that CoR-Fuse provides competitive performance against both generic fusion strategies and task-oriented fusion frameworks.
We additionally compare with ELF~\citep{luo2025elf} using its prescribed fixed encoder composition, \{H-optimus-0, Virchow2, UNI, GigaPath, CONCH-v1.5\}. Under this setting, CoR-Fuse achieves the highest AUC, while remaining competitive in AUPRC and F1 (Appendix Table~\ref{tab:elf_uvm}).

We fix the number of fused encoders to five and evaluate three representative encoder compositions spanning high, intermediate, and heterogeneous representation similarity, compositions were selected according to pre-defined CKA ranges before evaluating fusion performance (Appendix Table~\ref{tab:composition_summary}).
We also report the mean and best individual-FM AUC within each composition to characterize the predictive strength of the selected encoders.
For the high-similarity encoder composition \{H-optimus-0, H-optimus-1, UNI, GigaPath, Phikon-v2\}, attention-based fusion yields higher performance than simpler aggregation strategies (Appendix Table~\ref{tab:uvm_high_similarity}).
For the heterogeneous composition \{Hibou-L, Virchow2, CONCH-v1.5, UNI-v2, H-optimus-0\}, simpler fusion strategies outperform attention-based methods (Appendix Table~\ref{tab:uvm_high_diversity}), indicating that complex interaction modeling is not universally beneficial.
Finally, we evaluate an intermediate-similarity composition \{H-optimus-1, H-optimus-0, Virchow2, UNI-v2, CONCH-v1.5\} (Appendix Table~\ref{tab:uvm_n5_balanced_group}). As shown in Table~\ref{tab:uvm_composition_summary}, across all three representative compositions, CoR-Fuse achieves competitive performance with low additional computational overhead at the fusion stage.

\begin{table}[htbp]
\centering
\caption{Summary of fusion performance across three representative five-encoder compositions on TCGA-UVM.
Values are reported as mean $\pm$ standard deviation across the three compositions.}
\label{tab:uvm_composition_summary}
\begin{tabular}{lccccc}
\toprule
Method & AUC & AUPRC & F1 & Params(M) & FLOPs(G)\\
\midrule
Mean Fusion & $0.82 \pm 0.03$ & $0.86 \pm 0.03$ & $0.73 \pm 0.02$ &$4.01$ &$72.34$\\
Concatenation & $0.81 \pm 0.04$ & $0.86 \pm 0.03$ & $0.73 \pm 0.04$ &$5.32$ &$96.03$\\
Global Weight & $0.82 \pm 0.03$ & $0.86 \pm 0.04$ & $0.74 \pm 0.03$ &$4.01$ &$72.34$\\
Gated Fusion & $0.81 \pm 0.02$ & $0.85 \pm 0.01$ & $0.73 \pm 0.02$ &$4.01$ &$72.34$\\
Cross Attention & $0.83 \pm 0.03$ & $0.86 \pm 0.02$ & $0.74 \pm 0.06$ &$6.11$ &$148.55$\\
Self Attention & $0.81 \pm 0.04$ & $0.85 \pm 0.03$ & $0.73 \pm 0.03$ &$24.19$ &$2235.86$\\
CoR-Fuse & $\mathbf{0.83 \pm 0.01}$ & $\mathbf{0.87 \pm 0.01}$ & $\mathbf{0.75 \pm 0.02}$ &$4.07$ &$72.34$\\
\bottomrule
\end{tabular}
\end{table}

\textbf{Fusion Performance with Different Foundation Model Set Sizes.}
\label{scaling law section}
We systematically evaluate fusion strategies under different foundation-model set sizes by increasing the number of integrated encoders from two to ten, following the ranking of individual encoder performance on the validation set.
As shown in Appendix Figure~\ref{fig:scaling_law}, increasing the number of foundation models does not necessarily improve fusion performance, and the benefit of additional encoders depends on the encoder composition and fusion strategy.
Across encoder-set sizes, CoR-Fuse exhibits the smallest empirical AUC dispersion among the evaluated methods, suggesting more consistent performance across different FM compositions (Appendix Table~\ref{tab:scaling_law_variance}).

\begin{table}[htbp]
\centering
\caption{Fusion performance under a representative five-encoder composition on TCGA-BLCA.}
\label{tab:blca_n5}
\begin{tabular}{lccccc}
\toprule
Method & AUC & AUPRC & F1 & Params(M) & FLOPs(G)\\
\midrule
Mean Fusion & $0.74 \pm 0.06$ & $0.74 \pm 0.10$ & $0.60 \pm 0.10$ &$4.01$ &$74.81$\\
Concatenation & $0.74 \pm 0.06$ & $0.75 \pm 0.08$ & $0.57 \pm 0.23$ &$5.32$ &$99.31$\\
Global Weight & $0.74 \pm 0.05$ & $0.75 \pm 0.09$ & $0.65 \pm 0.07$ &$4.01$ &$74.81$\\
Gated Fusion & $0.73 \pm 0.06$ & $0.74 \pm 0.10$ & $\mathbf{0.68 \pm 0.06}$ &$4.01$ &$74.81$\\
Cross Attention & $0.74 \pm 0.03$ & $0.73 \pm 0.06$ & $0.57 \pm 0.17$ &$6.11$ &$153.62$\\
Meta-Encoder & $0.64 \pm 0.06$ & $0.63 \pm 0.16$ & $0.63 \pm 0.11$ &$24.19$ &$2235.86$\\
FuseCPath & $0.67 \pm 0.06$ & $0.67 \pm 0.16$ & $0.59 \pm 0.15$ &$7.16$ &$3.21$\\
CoR-Fuse & $\mathbf{0.75 \pm 0.07}$ & $\mathbf{0.76 \pm 0.12}$ & $0.68 \pm 0.08$ &$4.04$ &$74.81$\\
\bottomrule
\end{tabular}
\end{table}

\subsubsection{Evaluation on TCGA-BLCA Cohort}
To evaluate fusion under a practically motivated, limited yet diverse encoder composition, we select five representative pathology foundation models spanning different architectures, pretraining strategies, and model scales: \{UNI, Virchow2, Midnight-12k, GigaPath, Phikon\}.
More detailed information is provided in Appendix~\ref{sec:blca_results}.
As shown in Table~\ref{tab:blca_n5}, CoR-Fuse achieves competitive performance and lower computational cost than Cross Attention and Meta-Encoder.

\subsection{Ablation Study}
\begin{table}[htbp]
\centering
\caption{Ablation study of CoR-Fuse components.}
\label{tab:ablation}
\begin{tabular}{lcccc}
\toprule
\multirow{2}{*}{Method} & \multicolumn{2}{c}{TCGA-UVM} & \multicolumn{2}{c}{TCGA-BLCA} \\
\cmidrule(lr){2-3} \cmidrule(lr){4-5}
 & AUC & AUPRC  & AUC & AUPRC \\
\midrule
w/o Consensus & $0.79 \pm 0.22$ & $0.83 \pm 0.20$ & $0.72 \pm 0.03$ & $0.72 \pm 0.10$\\
w/o Warm Start & $0.73 \pm 0.23$ & $0.77 \pm 0.22$ & $0.72 \pm 0.02$ & $0.71 \pm 0.11$\\
CoR-Fuse & $\mathbf{0.81 \pm 0.09}$ & $\mathbf{0.84 \pm 0.10}$ & $\mathbf{0.73 \pm 0.03}$ & $\mathbf{0.72 \pm 0.09}$\\
\bottomrule
\end{tabular}
\end{table}

To evaluate the contribution of each design component in CoR-Fuse, we conduct ablation studies on TCGA-UVM and TCGA-BLCA using the encoder composition: \{UNI, Virchow2, Midnight-12k, GigaPath, Phikon\} (Table~\ref{tab:ablation}).

\textbf{Effect of Consensus Representation.}
Removing the consensus component leaves the overall concentration of slide-level routing weights largely unchanged, with the mean Gini coefficient changing only from (Appendix~\ref{sec:ablation_study}). 
Nevertheless, it induces a modest redistribution of routing mass across foundation models: removing consensus decreases the mean weights assigned to Phikon and UNI, while increasing those assigned to Virchow2 and Midnight-12k. 
Thus, consensus conditioning primarily affects encoder-level weight allocation, while having limited influence on the overall concentration of the routing distribution. 
Thus, the consensus component primarily affects encoder-level weight allocation rather than the overall concentration of the routing distribution.

\textbf{Effect of Task-Driven Warm Start.}
Although final performance was comparable, warm-start initialization substantially improved training reliability.
It reduced validation AUC variance by 47.9\% on average, yielding more stable and predictable optimization trajectories.
More detailed discussion is provided in Appendix~\ref{sec:warm_start}. 

Overall, the ablation results suggest that consensus and warm-start initialization may play different roles: consensus appears to influence routing behavior, whereas warm start can facilitate optimization during training. Their effects on final predictive performance vary across the evaluated compositions and datasets.